\Askhsh{36571}{2}
\begin{ylh}{1}
    \item 2.3 Μέτρα Θέσης και Διασποράς
\end{ylh}
Οι βαθμοί των γραπτών δοκιμασιών στο μάθημα των μαθηματικών ενός μαθητή για το Α' τετράμηνο είναι $16,18,18,12,16,10$, και αποτελούν τις τιμές ενός δείγματος. Αν η μέση τιμή των βαθμών είναι $\overline{x} =15$ και η τιμή της διασποράς $s^2 = 9$.
\begin{alist}
    \item Να υπολογίσετε την τυπική απόκλιση $s$ και τον συντελεστή μεταβολής $CV$ του δείγματος των βαθμών. \monades{10}
    \item Να βρείτε τη διάμεσο $\delta$ του δείγματος. \monades{5}
    \item Οι βαθμοί των γραπτών δοκιμασιών στο μάθημα των μαθηματικών του μαθητή για το Β' τετράμηνο αυξήθηκαν όλοι κατά δύο μονάδες. Να υπολογίσετε τη νέα μέση τιμή των βαθμών του Β' τετράμηνου. \monades{10}
\end{alist}

